\documentclass[10pt,a4paper]{amsart}
\usepackage[margin=19mm]{geometry}
\usepackage{amsmath,amssymb,mathtools}
\usepackage[scr=rsfs,cal=euler]{mathalfa}
\usepackage{xfrac}
\usepackage[unicode,hidelinks,bookmarks=false]{hyperref}
\newcommand{\ie}{i.e.\ }
\newcommand{\cf}{cf.\ }
\newcommand{\resp}{respectively\ }
\newcommand{\Subsection}{\subsection}
\providecommand{\nobreakdash}{\nobreak\hbox{-}}
\setlength{\emergencystretch}{2em}
\multlinegap0pt
\usepackage{amssymb}
\usepackage{mathtools}
\usepackage[shortlabels]{enumitem}
\newtheorem{lemma}{Lemma}
\newcommand{\R}{\mathbb{R}}
\renewcommand{\epsilon}{\varepsilon}
\title{Classification of Steady Gradient Ricci-Yang-Mills Solitons on Surfaces}
\author{Michael Womack}
\date{Source: arXiv:2604.25041v2 (CC BY 4.0)}
\begin{document}
\maketitle

\noindent\emph{Source and licence.} Classification of Steady Gradient Ricci-Yang-Mills Solitons on Surfaces. Version arXiv:2604.25041v2 (CC BY 4.0), distributed by arXiv under the Creative Commons Attribution 4.0 International licence, http://creativecommons.org/licenses/by/4.0/, as recorded in the arXiv licence metadata for this exact version and shown on the abstract page https://arxiv.org/abs/2604.25041v2. Full licence notice: \texttt{licenses/arxiv-2604.25041-CC-BY-4.0.txt}.\par\smallskip

\noindent\textit{Editorial scope.} Counted unit: the article's lemma f\_crit\_point\_pos\_lemma (source0.tex lines 881--887). The article's complete original proof of it is delivered with it (source0.tex lines 888--930, one \texttt{proof} environment of 43 lines). No mathematics is written, repaired or re-derived: the statement and the proof are the article's own lines, in the article's own notation, and no retained line is substituted or re-flowed.  No further source-local context is needed: every cross-reference of every retained line resolves inside the two delivered spans.  Nothing was omitted from any retained span, so the package carries no omission record.  No retained line contains a citation, so the wrapper carries no bibliography and no bibliography entry.  No retained line contains a figure environment, an image-inclusion command or drawing code, so no image is delivered and the package declares no asset dependency.  Editorial formatting only, in addition: an AMS \texttt{amsart} preamble that restates the article's own declarations and macros used by the retained lines, namely the environments \texttt{lemma} on separate counters, together with the standard packages those lines need.  The retained environments print in this document as Lemma 1 in order of appearance, whereas the article prints the same items with the numbering of its own full document.  The complete native source of the article as distributed in the arXiv e-print for this exact version is bundled unchanged as \texttt{sources/199267/source0.tex}.
\begin{lemma}[$B>0$ Case]\label{f_crit_point_pos_lemma}
  Let $f$ be a solution to
  \[
    f'' = B-Ae^{2f}+\frac{1}{2}(f')^{2}
  \]
  where $B>0$. Then $f$ must have a critical point.
\end{lemma}

\begin{proof}
  Without loss of generality, we can assume that we specify initial conditions at $r=0$, i.e. $f(0)=c_1$ and $f'(0)=c_2$. Note that $f(-r)$ also solves the equation with initial conditions $f(0)=c_0$ and $f'(0)=-c_1$, which means we can reflect solutions to flip the sign of the derivative. 

  When $f>\tfrac{1}{2}\ln(\tfrac{B}{A}):=f_{th}$ this implies $B-Ae^{2f}<0$ (since $B-Ae^{2f_{th}}=0$). We will refer to $f_{th}$ as the \textit{threshold}. If for some $r\in\R$ we have $f(r)<f_{th}$ and $f'(r)<0$, this implies $f''(s)$ is uniformly bounded from below by a positive constant for $s>r$ until $f$ hits a critical point. The condition that $f'(r)<0$ can always be achieved by reflecting $f$ if necessary. Thus in this case, $f'$ must reach zero in finite time, and we are guaranteed a critical point in this case.

  We now consider the case that $f(0)>f_{th}$ and $f''(0)<0$. Reflecting the solution does not change $f''(0)$, so we can take $f'(0)<0$. In this case $f$ will decrease as $r$ increases, and if $f''\leq 0$ is preserved then $f$ must decrease below the threshold eventually which would imply $f$ is now in the first situation above. In fact, the only obstruction to reaching a critical point must be the case where $f(r)>f_{th}$, $f'(r)<0$ for all $r>0$ with $\lim_{r\to\infty}f(r)=f_{th}$ and $\lim_{r\to\infty}f'(r)=0$ (which implies $\lim_{r\to\infty}f''(r)=0$). If $f''$ is positive for too long, then $f$ clearly hits a critical point, and if $f''$ possibly oscillates about 0 (but does not limit to zero), either $f$ hits a critical point or $f$ decreases below the threshold (which then implies $f$ will hit a critical point).

  Thus it is sufficient to prove that in the setting of $f(r)>f_{th}$, $f'(r)<0$ for all $r$, $\lim_{r\to\infty}f(r)=f_{th}$ and $\lim_{r\to\infty}f'(r)=\lim_{r\to\infty}f''(r)=0$ that $f$ must reach a critical point. Suppose now for the sake of contradiction that we are in this case but $f$ does not have a critical point. By assumption there exists a sequence $\{r_{i}\}$ such that $\lim_{i\to\infty}f''(r_{i})=0$ and $\lim_{i\to\infty}f'(r_{i})=0$. Let $\epsilon>0$, and we can find some $r_{k}$ such that $0<f''(r_{k})<\epsilon$ and $-\epsilon<f'(r_{k})<0$. 

  We start by using the setup to determine stricter bounds on $f$ and $f'$ from the ODE in this setting. Since $f''(r_{k})>0$ we have that
  \[
    -2(B-Ae^{2f})<(f')^{2} \implies f'<-\sqrt{-2(B-Ae^{2f})}.
  \]
  Applying the bound on $f'$ yields a bound on $f$: 
  \[
    -\epsilon<f'<0 \implies -\epsilon<-\sqrt{-2(B-Ae^{2f})}
  \]
  which we solve for $f$ to yield
  \[
    f<\tfrac{1}{2}\ln\left(\tfrac{1}{A}\left(\tfrac{\epsilon^{2}}{2}+B\right)\right).
  \]
  Note that these inequalities rely on $f''>0$. Since $f^{(3)}=f'(f''-2Ae^{2f})$, the assumption that $f'<0$ along with the bounds at $r_{k}$ yield that $f^{(3)}(r_k)>0$. In order for $f^{(3)}$ to become negative (and before $f$ has a critical point), it must be that $f''-2Ae^{2f}>0$, and in particular this requires $f''>2B$ (since $f>f_{th}$). This is not possible under the assumption that $|f'|<\epsilon$ for $\epsilon$ sufficiently small, so $f''>0$ must hold in the regime we consider here.

  The goal now is to show $f'$ is sufficiently large that it forces $f$ below the threshold in finite time. From the bounds above, the distance from $f$ to the threshold value is less than
  \[
    \tfrac{1}{2}\ln\left(\tfrac{1}{A}\left(\tfrac{\epsilon^{2}}{2}+B\right)\right) - \tfrac{1}{2}\ln\left(\tfrac{B}{A}\right) = \tfrac{1}{2}\ln\left(\tfrac{\epsilon^{2}}{2B}+1\right) := \delta(\epsilon, B)
  \]
  so if $f$ does not move below the threshold, it must be that
  \[
    f(r)-f(r_k) = \int_{r_{k}}^{r}f'(t)dt > -\delta
  \]
  for all $r>r_{k}$. Break the interval $[f_{th},f_{th}+\delta]$ into subintervals $[f_{th}+\tfrac{\delta}{2^i}, f_{th}+\tfrac{\delta}{2^{{i-1}}}]$. The bound on $f'$ can be used to bound the time it takes to drop $f$ through these subintervals by a convergent geometric series, which implies $f$ must go below the threshold in finite time. Define $r_{i}$ to be the time that $f(r_{i})=f_{th}+\tfrac{\delta}{2^i}$. Recalling that $f''>0$, we have that on the $i$-th subinterval,
  \begin{align*}
    -\frac{\delta}{2^{i}} = \int_{r_{i-1}}^{r_{i}}f'(t)dt &< \int_{r_{i-1}}^{r_{i}}-\sqrt{-2(B-Ae^{2f})}\;dt\\
      &< (r_{i}-r_{i-1})\left(-\sqrt{-2(B-Ae^{2(f_{th}+\delta/2^{i})})}\right)\\
      &= (r_{i}-r_{i-1})\left(-\sqrt{-2B(1-e^{\delta/2^{i}})}\right)   \tag*{(using $Ae^{2f_{th}}=B$)}
  \end{align*}
  which implies 
  \[
    r_{i}-r_{i-1}< \frac{\delta}{2^{i}}\frac{1}{\sqrt{2B}}\frac{1}{\sqrt{e^{\delta/2^{i}}-1}}<\frac{\delta}{2^{i}}\frac{1}{\sqrt{2B}}\frac{1}{\sqrt{\delta/2^{i}}}=\sqrt{\frac{\delta}{2B}}\left(\frac{1}{\sqrt{2}}\right)^{i}.
  \]
  Summing these and noting that the right side is the desired convergent geometric series implies that there is a sequence of points $\{r_{i}\}$ such that $\lim_{i\to\infty}f(r_{i})=f_{th}$ and $\lim_{i\to\infty}r_{i}=R<\infty$. This contradicts the assumption that $f$ remains above the threshold. Since negating any of the assumptions that set up the contradiction implies the existence of a critical point, we are done.
\end{proof}

\end{document}
